\chapter{Course Overview and Data Analysis Workflow}

Acquiring neural data is only the first step. Once an experiment has produced spike trains, local field potentials, calcium traces, or behavior-aligned recordings, the real work begins: organizing the data, defining the response variables, choosing sensible summary statistics, and building analyses that survive contact with noise and biological variability.

The central workflow is not ``collect data, then run a model.'' It is a sequence of modeling decisions: what counts as a response, which time origin matters, which variability is signal rather than nuisance, and which assumptions justify the statistics being reported.

\begin{figure}[h]
\centering
\begin{tikzpicture}[
  node distance=1.1cm,
  workflowstep/.style={draw, rounded corners, align=center, minimum width=2.8cm, minimum height=1cm},
  arrow/.style={-{Latex[length=2mm]}, thick}
]
\node[workflowstep] (raw) {Raw recording\\spikes, LFP, behavior};
\node[workflowstep, right=of raw] (align) {Alignment\\events and trials};
\node[workflowstep, right=of align] (repr) {Response\\representation};
\node[workflowstep, below=of repr] (stats) {Statistics\\rates, variability, correlations};
\node[workflowstep, left=of stats] (model) {Model\\encoding or decoding};
\node[workflowstep, left=of model] (check) {Validation\\residuals and controls};
\draw[arrow] (raw) -- (align);
\draw[arrow] (align) -- (repr);
\draw[arrow] (repr) -- (stats);
\draw[arrow] (stats) -- (model);
\draw[arrow] (model) -- (check);
\draw[arrow] (check) -- (raw);
\end{tikzpicture}
\caption{A disciplined neural-data workflow is cyclic. Validation often sends the analysis back to earlier choices about alignment, representation, or summary statistics.}
\end{figure}

\section{Analysis philosophy}

The guiding principle is that prose and computation should reinforce each other. A derivation should make clear what an estimator means, and code should make the estimator's behavior visible on simulated or real data. When a preprocessing choice changes the scientific interpretation, the notes should say so explicitly rather than hiding the choice inside a script.

\begin{aside}[Core workflow]
Start from raw measurements, define the neural response representation, build summary statistics, and only then fit heavier models. Clean abstractions early make later inference much easier.
\end{aside}

\section{Canonical analysis loop}

A useful first-pass workflow for many neural datasets is:
\begin{enumerate}
  \item load and validate raw data,
  \item align signals to task-relevant events,
  \item compute simple descriptive statistics,
  \item visualize structure before fitting models,
  \item and keep the full pipeline reproducible.
\end{enumerate}

This sequence sounds obvious, but it guards against many common mistakes. It is often more informative to examine rasters, marginal distributions, and trial averages before building a large decoding model. Classic computational-neuroscience texts repeatedly emphasize that representation choices shape every downstream conclusion \parencite{dayan2001theoretical,gerstner2014neuronal,paninski2007statistical}.

\begin{keyequation}[Canonical rate estimate]
\[
r(t) = \sum_{i=1}^{n} w(t-t_i)
\]
\tcblower
Here $w$ is a smoothing kernel and $t_i$ are spike times. This simple construction turns a point process into a continuous signal that can be visualized, compared across conditions, and passed into downstream models.
\end{keyequation}

\section{Reusable computation}

Reusable code should live at the level of concepts, not one-off figures. Natural expansion points are dataset loaders, event-alignment helpers, spike-train feature extraction, plotting utilities for rasters and correlograms, and reproducible studies that document each analysis decision. The goal is not to automate interpretation away; it is to make the mechanical parts reliable enough that the scientific assumptions are exposed.

\section{Next steps}

As the notes grow, each lecture block should end with three things: the mathematical object introduced, the empirical diagnostic that estimates it, and the failure mode to check before trusting it. That convention keeps the course focused on analysis rather than notation alone.
